\documentclass[conference]{IEEEtran}
\usepackage{amsmath}
\usepackage{graphicx}
\usepackage{booktabs}
\usepackage{cite}
\usepackage{url}
\usepackage[hidelinks]{hyperref}

\begin{document}

\title{The Memory-Constrained Compressed Symbol Table\\for Embedded Compilers\\
\large BUDGET-SYM: An Adaptive Scope-Aware Symbol Representation Framework}

\author{
\IEEEauthorblockN{Shubhi Singh}
\IEEEauthorblockA{SCOPE, VIT Vellore\\Vellore, India}
\and
\IEEEauthorblockN{Adhyan Jain}
\IEEEauthorblockA{SCOPE, VIT Vellore\\Vellore, India}
\and
\IEEEauthorblockN{Akshith Venkataramana}
\IEEEauthorblockA{SCOPE, VIT Vellore\\Vellore, India}
\and
\IEEEauthorblockN{Arjuman}
\IEEEauthorblockA{SCOPE, VIT Vellore\\Vellore, India}
}

\maketitle

\begin{abstract}
Symbol tables are a core compiler data structure, but their memory footprint
can become significant in resource-constrained embedded compilation
environments, where conventional string-heavy storage competes with a
compiler's limited memory budget. This paper documents BUDGET-SYM, a
prototype adaptive, scope-aware symbol-table architecture that selects a
per-symbol representation---INLINE, INTERNED, or COMPRESSED---using memory
pressure, scope/lifetime information, and access frequency, including
post-insertion promotion of frequently accessed (``hot'') symbols. The
prototype is implemented in header-only C++14 and evaluated against two real
baselines across eight benchmark identifier workloads and a four-variant
ablation study. Using a documented, deterministic tracked-memory accounting
model (explicitly not OS resident-set size), the prototype reduces tracked
memory by 1.24$\times$--2.59$\times$ relative to the conventional baseline.
The ablation study indicates that scope-based reclamation is the dominant
contributor to memory savings. Review-2 extensions reported here include:
(1) a systematic grid-search sweep of all six policy thresholds producing
data-driven replacements for the original hand-picked constants; (2)
evaluation against real-world embedded C identifier corpora extracted from
FreeRTOS, the Arduino core library, and the Zephyr RTOS; and (3) multi-seed
statistical benchmarking with 95\% confidence intervals establishing
statistical significance. Review-2 (latency/ML extension) added two further
mechanisms: a 64-entry LRU lookup cache that measurably reduces BudgetSym's
own lookup latency on hot-path access patterns, and a workload-adaptive ML
threshold predictor (Ridge regression trained offline on 300 synthetic
workloads over a 5-feature online profile) that reconfigures the policy
thresholds after a 100-symbol profiling window with no additional runtime
dependency. This revision (Review-3) reports a correctness fix -- the
lookup cache did not account for scope shadowing, so a symbol shadowing an
outer-scope declaration of the same name could resolve to the wrong, stale
entry, confirmed empirically and fixed with a permanent regression test --
and a five-way comparison of alternative index hash functions and storage
structures (two additional hash functions, a Robin Hood open-addressing
baseline, and a shared-trie baseline) directly answering Review-2's own
proposed future work on reducing cold-lookup cost structurally. No
algorithm in that comparison wins both memory and lookup latency on any of
the eight datasets: Robin Hood wins lookup latency universally by forgoing
compression, and BudgetSym wins compression universally by paying the
reconstruction cost Robin Hood avoids. Lookup latency overhead remains an
honestly reported cost of the design, now examined from two independent
angles (caching and structural alternatives) rather than eliminated.
\end{abstract}

\begin{IEEEkeywords}
Symbol Table, Embedded Compilers, Memory-Constrained Computing, Compiler Data
Structures, String Interning, Adaptive Data Structures, Symbol Compression,
Policy Threshold Optimization, Real-World Workload Evaluation, Lookup
Caching, Machine-Learned Configuration
\end{IEEEkeywords}

\section{Introduction}
A compiler repeatedly creates, resolves, and updates symbolic bindings while
translating a program. A symbol table stores the name of each identifier
together with attributes such as type, scope, declaration information, and
storage-related metadata. Classical compiler-construction guidance treats
memory space and lookup/insertion time as competing design criteria for
symbol tables \cite{aho}. This tension is amplified when compilation targets
or toolchains operate under tight memory constraints, as is common for
embedded systems \cite{gcc-msp430}.

A conventional hash table (e.g., \texttt{unordered\_map<string, Symbol>})
incurs overhead from buckets, allocator metadata, and one full string copy
per symbol. Alternative representations exploit different forms of
redundancy: string interning stores one copy of each distinct identifier
\cite{rust-interning}; prefix-based storage exploits shared prefixes between
identifiers \cite{liu-lab2}. Each technique individually addresses only part
of the problem: interning helps only on exact repeats, and most identifiers
in a translation unit are declared once; a uniformly applied compression
scheme does not ask whether a specific identifier's lookup cost is worth
paying right now.

This motivates an adaptive approach. BUDGET-SYM is a prototype that selects
a representation per symbol based on memory pressure, scope/lifetime,
exact-repeat status, and prefix similarity, and subsequently promotes
symbols to a faster representation if they turn out to be accessed
frequently. Review-1 established the prototype and baseline evaluation.
Review-2 addressed the three most significant limitations identified in
Review-1: hand-picked thresholds, synthetic-only workloads, and absence of
statistical significance testing, and added the lookup cache and ML
threshold predictor along the lookup-overhead axis Review-1 had left
unresolved. This revision (Review-3) reports a correctness fix found while
reviewing that cache (a scope-shadowing bug, confirmed empirically and
fixed with a permanent regression test) and directly answers Review-2's
own Future Work proposal to reduce cold-lookup cost structurally, by
implementing and measuring two alternative storage structures (Robin Hood
open addressing, a shared trie) and two additional index hash functions
against the same benchmark suite.

The main contributions of this work are:
\begin{enumerate}
\item A working prototype BUDGET-SYM, implemented in header-only C++14
  [Review-1].
\item A unified representation-selection policy spanning INLINE, INTERNED,
  and COMPRESSED representations [Review-1].
\item Scope/lifetime-aware symbol reclamation [Review-1].
\item A deterministic tracked-memory accounting model distinct from OS
  resident memory [Review-1].
\item Access-driven post-insertion promotion of frequently looked-up
  (``hot'') symbols [Review-1].
\item A four-variant ablation study isolating mechanism contributions
  [Review-1].
\item Systematic grid-search sweep of all six policy thresholds, replacing
  hand-picked constants with data-driven optima [Review-2].
\item Multi-seed statistical benchmarking with 95\% confidence intervals
  [Review-2].
\item Evaluation against real embedded C identifier corpora from FreeRTOS,
  Arduino core, and Zephyr RTOS [Review-2].
\item A 64-entry, hash-keyed LRU lookup cache that memoizes recently-resolved
  symbols, directly reducing the reconstruct-and-compare cost that dominates
  BUDGET-SYM's lookup latency [Review-2].
\item Per-entry memoized COMPRESSED-chain reconstruction, eliminating
  redundant chain walks on repeated lookups of the same compressed symbol
  [Review-2].
\item An online WorkloadProfiler (5 features, 100-symbol sampling window)
  feeding a Ridge-regression ThresholdPredictor trained offline on 300
  synthetic workloads, reconfiguring policy thresholds mid-compilation with
  no runtime ML dependency [Review-2].
\item A correctness fix for a scope-shadowing bug in the Review-2 lookup
  cache, found by review, confirmed empirically, fixed, and covered by a
  permanent regression test [\textbf{new}].
\item Two additional index hash functions (MurmurHash3, DJB2) as swappable
  BUDGET-SYM backends, and two new structural baselines (Robin Hood open
  addressing, a shared trie), benchmarked against the same eight datasets
  [\textbf{new}].
\end{enumerate}

\section{Background}
\subsection{Compiler Symbol Tables}
A symbol table binds each identifier encountered during compilation to
attributes required for semantic analysis and code generation, such as type,
scope, and storage information \cite{aho}. Design choices directly affect
both memory footprint and the speed of name resolution.

\subsection{Symbol Storage in Resource-Constrained Environments}
Embedded toolchains routinely expose memory-related compilation options and
operate under tighter memory budgets than desktop-scale compilation
\cite{gcc-msp430}. This provides a concrete setting in which symbol-table
overhead becomes a first-class engineering concern.

\subsection{String Interning}
String pools separate identifier text from symbol records so that identical
names are stored once and referenced by a compact handle
\cite{jones-notes,rust-interning}. The benefit is proportional to the
exact-duplication rate of identifiers in the workload.

\subsection{Prefix-Based / Compressed String Representations}
Tries and compressed tries (radix trees) exploit shared prefixes between
strings, trading pointer/traversal overhead for reduced storage on long,
sparse identifier sets \cite{liu-lab2}.

\subsection{Scope and Symbol Lifetime}
Lexical scoping means many symbols have a short, well-defined lifetime.
Scope-aware management can reclaim memory eagerly at scope exit rather than
holding every declared identifier for the duration of compilation.

\subsection{Memory--Latency Trade-offs}
Reducing memory per symbol typically increases the cost of reconstructing or
locating that symbol's data at lookup time. A central question for any
adaptive scheme is whether it can move the memory--latency trade-off curve
outward rather than simply sliding along it. This revision's lookup cache and
memoized reconstruction are a direct response to that question: rather than
changing the representation trade-off itself, they amortize the
reconstruction cost across repeated accesses to the same symbol.

\section{Literature Review and Related Work}
This review is based on established compiler-construction references and
documented industrial practice. Aho et al. \cite{aho} describe symbol-table
design as a trade-off between memory space, insertion/search time, and
scoping complexity, and describe the common design of a hash table paired
with a separate string pool. Course-level material \cite{liu-lab2,
jones-notes} documents this same baseline design. Interning is used in
production compilers: the Rust compiler represents symbols using compact
interned identifiers \cite{rust-interning}. Embedded toolchain documentation
\cite{gcc-msp430} confirms that memory-constrained compilation targets are a
current engineering concern.

\textbf{Research Gap.} The literature establishes the individual building
blocks in isolation. The narrower open question adopted here is whether a
compiler-specific, runtime/use-aware signal set---memory pressure,
scope/lifetime, exact-repeat status, prefix similarity, and access
frequency---can drive a unified per-symbol representation policy under an
explicit memory budget, while controlling lookup overhead, and whether that
control can be extended to a small, hash-keyed lookup cache and a
lightweight, offline-trained threshold predictor without adding a runtime ML
dependency. A fuller database-indexed prior-art search (IEEE Xplore, ACM DL,
arXiv) remains a planned step before any patentability claim.

\section{Problem Formulation}
Given a set of symbols $S$, a memory budget $B$, and a set of representation
choices $R = \{\text{INLINE}, \text{INTERNED}, \text{COMPRESSED}\}$, the
problem is to select a representation $r_i \in R$ for each symbol $s_i \in
S$ such that the total tracked memory does not exceed $B$, while minimizing
a combined cost:
\begin{equation}
C = \alpha M + \beta L + \gamma T
\end{equation}
where $M$ is memory cost, $L$ is lookup latency, $T$ is transition/re-encoding
overhead, and $\alpha, \beta, \gamma$ are configurable policy weights. The
current prototype uses a heuristic instantiation of Eq.~(1). Review-2
contributed a data-driven method for selecting the scalar thresholds that
control when each branch of the heuristic fires. This revision addresses the
$L$ term directly (via the lookup cache) and extends the threshold-selection
method from a one-time offline sweep to a per-workload, online-triggered
prediction.

\section{Proposed BUDGET-SYM Architecture}
The overall pipeline is: source identifiers are extracted and passed to a
central controller, which consults six subsystems before an Adaptive
Representation Policy commits each symbol to INLINE, INTERNED, or
COMPRESSED storage behind a compact index. Fig.~1 (see
\texttt{docs/architecture.md} in the accompanying repository) traces this
data flow. Each subsystem is described below as one paragraph.

\textbf{1) Memory Monitor.} Tracks budget, current tracked-memory usage, and
derived pressure (\texttt{MemoryTracker}, \texttt{include/memory\_tracker.hpp}).
Pressure is a deterministic ratio of current to budget, clamped to $[0,1]$,
and is the sole memory-side signal \texttt{decide()} consults.

\textbf{2) Scope/Lifetime Analyzer.} Maintains the scope stack
(\texttt{enterScope()}/\texttt{exitScope()}) and drives reclamation:
on scope exit, every symbol declared in that scope is tombstoned and, where
physically possible, its bytes are reclaimed from the tracker immediately.
This is the mechanism the Review-1 ablation study identified as the dominant
contributor to memory savings.

\textbf{3) Access Profiler.} Tracks a per-symbol \texttt{accessCount},
incremented on every successful \texttt{lookup()}. This is the signal that
drives post-insertion promotion from COMPRESSED to INTERNED once a symbol
crosses \texttt{hotAccessThreshold}.

\textbf{4) WorkloadProfiler + ThresholdPredictor (new).} An online,
constant-memory accumulator (\texttt{include/workload\_profiler.hpp})
observes the first 100 inserted identifiers and extracts five features: mean
name length, prefix-similarity coefficient, repeat rate, scope-depth
variance, and character-entropy. Once the window fills, these features are
passed to \texttt{ThresholdPredictor::predict()}
(\texttt{include/predicted\_thresholds.hpp}), a set of six hardcoded
Ridge-regression models trained offline (Section VIII-F), and the resulting
\texttt{PolicyConfig} replaces the running configuration for the remainder
of the compilation. No machine-learning library is linked into the compiler
itself; the predictor is six clamped linear expressions.

\textbf{5) LRU Lookup Cache (new).} A 64-entry, hash-and-name-keyed LRU cache
(\texttt{include/lookup\_cache.hpp}) sits in front of the normal
hash-bucket-then-reconstruct lookup path. A hit returns the cached entry
pointer directly, skipping both the bucket walk and, for COMPRESSED
symbols, the chain-reconstruction cost entirely. Entries are invalidated by
hash on scope exit, so a cache hit can never resurrect a reclaimed symbol.

\textbf{6) Adaptive Representation Policy (\texttt{decide()}).} Consumes
identifier length, exact-repeat status, prefix similarity to the previously
inserted identifier, current memory pressure, and the insert-time access-
frequency hint to choose INLINE, INTERNED, or COMPRESSED for each new
symbol, as described in Section VI.

\section{Adaptive Representation Policy}
Representation decisions are influenced by identifier length, exact-repeat
status, prefix similarity to already-stored identifiers, current memory
pressure, scope depth, and post-insertion measured access frequency.
\begin{itemize}
\item \textbf{INLINE} --- used when a small/cheap identifier does not
  justify the bookkeeping cost of interning or compression.
\item \textbf{INTERNED} --- used when exact identifier reuse makes shared,
  reference-counted storage beneficial.
\item \textbf{COMPRESSED} --- used when a long or prefix-similar identifier
  offers sufficient savings under front-coded/prefix-shared storage.
\end{itemize}
Post-insertion promotion moves a symbol from COMPRESSED to INTERNED once its
measured access count crosses the hot-access threshold, trading its memory
saving for faster lookup.

\section{Memory Accounting Model}
Memory is reported using a documented, deterministic tracked-memory model,
computed as the sum of string payload, symbol metadata, index storage,
representation-specific storage, and fixed metadata constants, consistent
across all three implementations. This is explicitly not a measurement of
full operating-system resident-set size (RSS); it is a controlled accounting
model chosen so that all three implementations are compared on a
like-for-like basis. This limitation is stated plainly rather than left
implicit. The lookup cache and per-entry memoization added in this revision
do not change this model: they add a small, fixed number of bytes per cached
or memoized entry, which is not currently charged against the tracked-memory
total (see Threats to Validity, item 8).

\section{Implementation}
The prototype is implemented in header-only C++14. It provides a common
symbol-table interface (\texttt{insert}, \texttt{lookup}, \texttt{remove},
\texttt{contains}, \texttt{enterScope}, \texttt{exitScope},
\texttt{memoryUsage}, \texttt{statistics}) implemented by the conventional
baseline, the interned baseline, and BUDGET-SYM itself, so that all three
are exercised with identical operation sequences. Development surfaced
several real toolchain/logic bugs, caught and fixed during implementation
--- including, in this revision, a latent measurement bug in the shared
benchmark harness itself (Section VIII-E).

\subsection{Review-2 Extension: Systematic Threshold Optimization}
\textbf{Review-1 Limitation.} The six policy thresholds in
\texttt{PolicyConfig}---\texttt{inlineMaxLen}, \texttt{compressMinLen},
\texttt{lowPressureThreshold}, \texttt{highPressureThreshold},
\texttt{hotAccessThreshold}, \texttt{prefixSimilarityMinShared}---were chosen
by inspection during development with no systematic justification.

\textbf{Approach.} A grid search is conducted over the six-dimensional
threshold space using the eight synthetic benchmark datasets as the
optimization target. Each configuration is scored by mean compression ratio
across all eight datasets. Total configurations evaluated: $6 \times 5
\times 4 \times 5 \times 5 \times 5 = 15{,}000$, shown in Table~\ref{tab:grid}.

\begin{table}[htbp]
\caption{Threshold Grid-Search Space}
\label{tab:grid}
\centering
\begin{tabular}{lll}
\toprule
Threshold & R1 Value & Sweep Range \\
\midrule
inlineMaxLen & 12 & 6, 8, 10, 12, 14, 16 \\
compressMinLen & 10 & 6, 8, 10, 12, 14 \\
lowPressureThreshold & 0.50 & 0.30, 0.40, 0.50, 0.60 \\
highPressureThreshold & 0.85 & 0.70, 0.75, 0.80, 0.85, 0.90 \\
hotAccessThreshold & 3 & 2, 3, 4, 5, 6 \\
prefixSimilarityMinShared & 4 & 2, 3, 4, 5, 6 \\
\bottomrule
\end{tabular}
\end{table}

\textbf{Results.} The grid-search-optimal configuration achieves a mean
compression ratio of 1.87$\times$ across the eight datasets, compared with
1.71$\times$ for the Review-1 hand-picked thresholds --- an improvement of
approximately 9\% without any change to the underlying architecture.
Optimized values are reported in \texttt{results/optimal\_policy.csv} in the
accompanying repository.

\textbf{Caveat.} The optimized thresholds are tuned to the eight synthetic
datasets. Generalization to real-world workloads is treated as an open
empirical question, evaluated in the next subsection. This revision's
threshold predictor (Section VIII-F) is a per-workload alternative to this
one-time, dataset-fixed sweep.

\subsection{Review-2 Extension: Real-World Workload Evaluation}
\textbf{Review-1 Limitation.} All eight benchmark datasets in Review-1 were
synthetically generated. Results on real embedded-compiler identifier
traffic were explicitly deferred.

\textbf{Corpora.} Three open-source embedded C codebases were vendored
locally under \texttt{corpora/} (see \texttt{docs/corpus\_setup.md}):
FreeRTOS-Kernel (655 files), the Arduino core (AVR) library (311 files), and
the Zephyr RTOS \texttt{kernel/} and \texttt{drivers/} subtrees (4{,}298
files). A Python script (\texttt{scripts/extract\_identifiers.py})
traverses \texttt{.c}/\texttt{.h} files, applies a regex tokenizer, filters
C89/C99 language keywords only (standard-library and platform typedefs such
as \texttt{uint32\_t} are deliberately kept, since compressing those is
exactly what BudgetSym is meant to do), and feeds the resulting sequence
--- source order, duplicates preserved --- into the same benchmark harness
used for the synthetic datasets (\texttt{corpus\_bench.exe}).

\begin{table}[htbp]
\caption{Real-World Corpus Results (Single Run, 64\,MB Budget)}
\label{tab:corpus}
\centering
\small
\begin{tabular}{lccc}
\toprule
Codebase & Tokens & Interned & BudgetSym \\
\midrule
FreeRTOS & 1{,}107{,}471 & 2.522$\times$ & 2.520$\times$ \\
Arduino (AVR) & 329{,}480 & 2.335$\times$ & 2.346$\times$ \\
Zephyr (kernel+drivers) & 3{,}710{,}619 & 2.315$\times$ & 2.312$\times$ \\
\bottomrule
\end{tabular}
\end{table}

\textbf{Findings.} Both plain interning and BudgetSym reduce tracked memory
by roughly 2.3$\times$--2.5$\times$ versus the conventional baseline on all
three real corpora --- a substantially larger absolute reduction than any of
the eight synthetic datasets show for interning alone (which sits at
$\sim$1.05$\times$, Table~\ref{tab:multiseed}). This is a direct consequence
of real embedded identifier streams' very high repeat rate: the same
handful of identifiers (loop counters, HAL field names, common helpers) are
declared and referenced repeatedly across hundreds of files, so a large
fraction of insertions are exact-repeat lookups that both Interned and
BudgetSym resolve identically (\texttt{decide()}'s first rule routes every
exact repeat to INTERNED, matching the Interned baseline's own behavior).

\textbf{Honest reading.} On these three corpora, BudgetSym's advantage over
plain interning is within noise ($\pm$0.03$\times$, and Zephyr's promotion
count was 0 in this single run) rather than the larger margins seen on the
synthetic datasets. BudgetSym's real benefit here is not differentiating
itself from interning on already-repeated identifiers --- it is scope-aware
reclamation and the COMPRESSED representation's handling of long,
prefix-similar \emph{first-occurrence} names, neither of which a flat
interning pool provides. The synthetic high-prefix-similarity and
nested-scopes datasets, not these three general-purpose corpora, remain the
workloads that isolate that advantage most clearly. This is a corrected
finding versus an earlier, unverified draft of this subsection, which
reported per-corpus BudgetSym-only ratios of
1.51$\times$/1.44$\times$/1.63$\times$ without an Interned baseline for
direct comparison; those figures could not be reproduced when the corpora
were actually run and are superseded by the numbers above.

\subsection{Review-2 Extension: Multi-Seed Statistical Validation}
\textbf{Review-1 Limitation.} All benchmark results in Review-1 came from a
single deterministic run. No confidence intervals or significance tests were
reported.

\textbf{Method.} Each synthetic benchmark dataset is re-generated with 30
independent random seeds. Mean and 95\% confidence intervals are computed
using the $t$-distribution ($n=30$).

\textbf{Findings.} BudgetSym's compression-ratio advantage over both
baselines is statistically significant ($p<0.01$) on all eight datasets
(Table~\ref{tab:multiseed}). The smallest advantage (random-identifiers,
$\sim$1.40$\times$) has a tight 95\% CI, confirming the result is not
attributable to a single favorable seed.

\begin{table}[htbp]
\caption{Multi-Seed Mean Compression Ratio vs. Conventional (30 Seeds, 95\% CI)}
\label{tab:multiseed}
\centering
\small
\begin{tabular}{lccc}
\toprule
Dataset & Interned & BudgetSym & $p$ vs. Conv. \\
\midrule
small & 1.058$\times$ & 1.810$\times$ & $<10^{-30}$ \\
medium & 1.056$\times$ & 1.783$\times$ & $<10^{-50}$ \\
large & 1.056$\times$ & 1.806$\times$ & $<10^{-70}$ \\
high-prefix-sim. & 1.045$\times$ & 2.595$\times$ & $<10^{-80}$ \\
random-identifiers & 1.054$\times$ & 1.403$\times$ & $<10^{-45}$ \\
nested-scopes & 1.053$\times$ & 2.042$\times$ & $<10^{-75}$ \\
hot-cold-access & 1.053$\times$ & 1.763$\times$ & $<10^{-50}$ \\
memory-stress & 1.048$\times$ & 1.586$\times$ & $<10^{-75}$ \\
\bottomrule
\end{tabular}
\end{table}

Note: these Review-2 multi-seed figures are produced by the current
codebase, in which BudgetSym's post-100-symbol thresholds are already
reconfigured by this revision's threshold predictor (Section VIII-F) unless
explicitly held fixed; they are consequently slightly higher than the
Review-2 draft's original hand-picked-only numbers and should be read as
current-system, not purely-Review-2-policy, figures.

\subsection{Review-2 Extension: Latency Reduction via Lookup Cache}
\textbf{Motivation.} Every prior revision of this paper reported lookup
latency as an unresolved, honestly-stated cost: BUDGET-SYM's
hash-then-reconstruct lookup path is structurally slower than a plain hash
table for any non-INLINE representation, because a COMPRESSED symbol must
walk its front-coding chain to confirm a match (Section I of
\texttt{docs/architecture.md}). This subsection reports the first concrete
mitigation.

\textbf{Design.} A 64-entry LRU cache (\texttt{LRULookupCache<64>},
\texttt{include/lookup\_cache.hpp}) is keyed by the 64-bit FNV-1a hash of the
identifier plus the identifier itself, to guard against the non-cryptographic
hash's collision risk producing a false hit. A hit moves the entry to the
front of an internal \texttt{std::list} and returns its (stable, id-packed)
entry reference directly, bypassing both the hash-bucket walk and any
COMPRESSED-chain reconstruction. \texttt{exitScope()} invalidates by hash
every symbol it reclaims, so a stale cache entry can never resurrect a
reclaimed symbol. A companion, independent optimization---per-entry memoized
COMPRESSED reconstruction (\texttt{CompEntry::reconstructedCache})---caches
each chain node's fully-reconstructed string the first time it is computed,
so repeated lookups of the same compressed symbol, or lookups that
transitively re-walk a shared chain prefix, no longer re-walk the chain at
all.

\textbf{Measurement bug found and fixed.} While instrumenting this
revision's benchmark pass, a latent bug was found in the existing shared
benchmark harness (\texttt{include/bench\_metrics.hpp}): the lookup-timing
loop was written as \texttt{sinkHit = sinkHit || table.lookup(id)}, which
short-circuits as soon as the first lookup succeeds, silently skipping every
subsequent \texttt{lookup()} call in the timed region for the rest of that
loop. This means every previously reported \texttt{lookup\_success\_us}
figure, across all revisions of this project, measured an almost-empty loop
rather than repeated lookup work. It is fixed in this revision (evaluated
correctly for the first time here); prior lookup-latency numbers in Review-1
and Review-2 should be treated as unreliable and are not re-cited in this
revision's own tables.

\textbf{Results.} Table~\ref{tab:cache} reports per-dataset lookup latency
with the cache enabled (BudgetSym-WithCache) against the cache disabled
(BudgetSym) and the conventional baseline, using a hot-path access pattern
(32 symbols, 20 repeated rounds --- see \texttt{results/benchmark\_results\_v3.csv}).
Across the eight synthetic datasets, the cache achieved a hit rate of 95\%
under this access pattern (one necessary cold miss per hot symbol, then all
hits) and reduced mean per-lookup latency relative to the same table with
the cache disabled on most datasets, though it did not close the gap to the
conventional baseline's uncompressed \texttt{unordered\_map} lookup on any
dataset -- consistent with the cache amortizing repeated cost rather than
eliminating the underlying reconstruct-and-compare structure.

\begin{table}[htbp]
\caption{Lookup Latency ($\mu$s/op, Hot-Path Pattern) --- With vs. Without Cache}
\label{tab:cache}
\centering
\small
\begin{tabular}{lccc}
\toprule
Dataset & Conventional & BudgetSym & +Cache (hit\%) \\
\midrule
small & 0.032 & 0.058 & 0.068 (95\%) \\
medium & 0.040 & 0.074 & 0.094 (95\%) \\
large & 0.049 & 0.089 & 0.078 (95\%) \\
high-prefix-sim. & 0.041 & 0.161 & 0.108 (95\%) \\
random-identifiers & 0.034 & 0.089 & 0.095 (95\%) \\
nested-scopes & 0.041 & 0.144 & 0.118 (95\%) \\
hot-cold-access & 0.040 & 0.173 & 0.139 (95\%) \\
memory-stress & 0.037 & 0.110 & 0.087 (95\%) \\
\bottomrule
\end{tabular}
\end{table}

\textbf{Honest reading.} The cache's effect splits along dataset character:
it reduces BudgetSym's own lookup latency by 12--33\% on five of the eight
datasets (large, high-prefix-similarity, nested-scopes, hot-cold-access,
memory-stress), with the largest reductions on the three datasets with
nonzero \texttt{promotions} in \texttt{results/benchmark\_results\_v3.csv}
(high-prefix-similarity, nested-scopes, hot-cold-access) -- i.e. where
COMPRESSED-chain activity, and hence chain-walk cost, is highest. It instead
\emph{increases} latency by roughly 6--26\% on small, medium, and
random-identifiers, whose sampled hot symbols are mostly INLINE/INTERNED
already cheap to reconstruct, so the cache's own bookkeeping (hash lookup,
list splice) costs more than the walk it replaces. This is a plausible,
dataset-dependent trade-off rather than pure noise, though with only a
single run per dataset here (see Threats to Validity, item 1, on multi-seed
validation still being owed for this claim), it should be read as a
direction, not a precise per-dataset percentage. It does not, on this measurement, bring
BudgetSym's lookup latency below the conventional baseline's on any dataset
--- that would require eliminating the reconstruct-and-compare step for cold
entries entirely, which is future work (Section XV).

\subsection{Review-2 Extension: ML-Driven Threshold Prediction}
\textbf{Motivation.} Section VIII-A's grid search produces one fixed,
dataset-averaged optimum. A workload whose characteristics differ from the
eight synthetic benchmark datasets gets no benefit from that sweep beyond
whatever generalizes coincidentally. This subsection replaces the one-time
sweep with a per-workload predictor triggered automatically during
compilation.

\textbf{WorkloadProfiler.} \texttt{include/workload\_profiler.hpp} observes
the first 100 inserted identifiers (matching this revision's sampling window)
and extracts five features in constant memory: mean name length,
prefix-similarity coefficient (fraction of consecutive insert pairs sharing
$\geq 4$ characters), repeat rate, scope-depth variance, and Shannon entropy
over the character distribution.

\textbf{Training procedure.} \texttt{scripts/train\_threshold\_predictor.py}
generates 300 synthetic workloads (100 identifiers each) with randomized
mean length ($[4,20]$), prefix-similarity coefficient ($[0,0.8]$), repeat
rate ($[0,0.5]$), scope-depth variance ($[0,5]$), and character entropy
target ($[2,5]$). For each workload, a Python port of BudgetSym's
\texttt{decide()}/cost model is grid-searched over a reduced,
144-combination threshold grid ($3\times3\times2\times2\times2\times2$,
coarser than Section VIII-A's 15,000-combination sweep for compute
tractability in an offline Python script) to find the compression-ratio-
optimal \texttt{PolicyConfig} for that workload. The resulting
(features, optimal-thresholds) pairs train one
\texttt{sklearn.linear\_model.Ridge(alpha=1.0)} model per threshold, with an
80/20 train/test split.

\textbf{Results.} Table~\ref{tab:r2} reports test-set $R^2$ per threshold.
Two thresholds (\texttt{lowPressureThreshold}, \texttt{hotAccessThreshold})
collapsed to a constant prediction: under the 64\,MB budget used for
training and the 100-symbol window, tracked memory never approaches either
pressure threshold, so the reduced grid search always picked the same value
for these two regardless of workload features --- an honest negative
finding, not a modeling failure, and consistent with Review-1's own note
that \texttt{hotAccessThreshold} was already near its optimum.
\texttt{compressMinLen} and \texttt{prefixSimilarityMinShared} show the
strongest fits, matching the intuition that these two thresholds are most
directly determined by the prefix-similarity and length features.

\begin{table}[htbp]
\caption{ThresholdPredictor Test-Set $R^2$ per Threshold (300 Workloads, 80/20 Split)}
\label{tab:r2}
\centering
\begin{tabular}{lc}
\toprule
Threshold & $R^2$ \\
\midrule
inlineMaxLen & 0.115 \\
compressMinLen & 0.243 \\
lowPressureThreshold & 0.000 (constant) \\
highPressureThreshold & 0.889 \\
hotAccessThreshold & 1.000 (constant) \\
prefixSimilarityMinShared & 0.197 \\
\bottomrule
\end{tabular}
\end{table}

Table~\ref{tab:modelcompare} compares all seven candidate models (PRD
Section 6.5 Step 2) on the same 300-workload, 80/20 split: mean $R^2$ and
mean absolute error across the six thresholds, and downstream compression
ratio when each model's predicted thresholds are applied to the eight real
benchmark datasets (profiled on each dataset's first 100 identifiers, same
as Section VIII-F's own methodology).

\begin{table}[htbp]
\caption{ML Model Comparison (Mean Across 6 Thresholds, 8-Dataset Downstream Compression)}
\label{tab:modelcompare}
\centering
\small
\begin{tabular}{lcccc}
\toprule
Model & Mean $R^2$ & Mean MAE & Downstream Ratio & Exportable \\
\midrule
Baseline Mean & 0.308 & 0.897 & 1.732$\times$ & Yes \\
Lasso & 0.337 & 0.831 & 1.753$\times$ & Yes \\
\textbf{Ridge} & \textbf{0.407} & \textbf{0.773} & \textbf{1.795$\times$} & \textbf{Yes} \\
Decision Tree & $-$48.98 & 0.427 & 1.714$\times$ & No \\
Random Forest & $-$1.917 & 0.439 & 1.724$\times$ & No \\
SVR & 0.482 & 0.645 & 1.750$\times$ & No \\
KNN & 0.651 & 0.606 & 1.719$\times$ & No \\
\bottomrule
\end{tabular}
\end{table}

\textbf{Reading the comparison.} Two non-linear models (SVR, $R^2=0.482$;
KNN, $R^2=0.651$) post higher raw $R^2$ than Ridge, but neither is
exportable as compile-time constants -- SVR needs its full support-vector
set and KNN needs the entire 300-workload training set at inference time,
both unacceptable for a header-only, no-runtime-dependency library. Among
the three exportable models, Ridge is the clear best on every column: it
beats Lasso's $R^2$ by 21\% relative and the baseline-mean predictor by
32\% relative, and it achieves the single highest downstream compression
ratio of \emph{any} model in the table, exportable or not (1.795$\times$
vs. KNN's 1.719$\times$) -- despite not having the best raw $R^2$. Decision
Tree and Random Forest both post strongly negative test-set $R^2$ (worse
than predicting the training mean), a genuine overfitting failure at this
training-set size (240 training rows after the 80/20 split) with these
un-tuned depths/tree-counts, not a bug in the comparison harness -- their
downstream compression ratios are unaffected because \texttt{compressMinLen}
and \texttt{prefixSimilarityMinShared}, the two thresholds that matter most
for compression ratio (Table~\ref{tab:r2}), still land in a serviceable
range even from an overfit model on this evaluation.

Table~\ref{tab:mlcompare} compares compression ratio on the
high-prefix-similarity dataset across four configurations: the conventional
1.00$\times$ reference, Review-1's hand-picked thresholds, Review-2's
dataset-averaged grid-search optimum, and this revision's ML-predicted
thresholds (profiled on the same dataset's first 100 identifiers, then
applied for the full run).

\begin{table}[htbp]
\caption{Compression Ratio, high-prefix-similarity Dataset}
\label{tab:mlcompare}
\centering
\begin{tabular}{lc}
\toprule
Configuration & Compression Ratio \\
\midrule
Conventional & 1.00$\times$ \\
Hand-Picked (R1) & 2.37$\times$ \\
Grid-Search (R2, dataset-averaged) & 2.59$\times$ \\
ML-Predicted (this revision) & 2.53$\times$ \\
\bottomrule
\end{tabular}
\end{table}

\textbf{Honest reading.} The ML-predicted configuration slightly
underperforms the dataset-specific grid-search optimum on this one dataset
(2.53$\times$ vs. 2.59$\times$), which is expected: the grid search's
15,000-combination sweep was optimized directly against this exact dataset,
while the predictor generalizes from 300 unrelated synthetic workloads
through a coarser 144-combination search and a linear model. Its practical
value is not beating a dataset-specific optimum but recovering most of that
optimum's benefit (2.53$\times$ vs. 2.37$\times$ hand-picked, a 6.8\%
improvement) automatically, at runtime, for a workload the grid search was
never run against, with no runtime ML dependency (\texttt{ThresholdPredictor::predict()}
is six clamped linear expressions).

\textbf{Feature ablation (PRD Section 6.5 Step 3).} To identify which of the
five \texttt{WorkloadFeatures} actually drive Ridge's predictions, each
feature is dropped in turn (retraining on the remaining four) and the mean
$R^2$ delta across all six thresholds is reported, along with which single
threshold's $R^2$ moved the most.

\begin{table}[htbp]
\caption{Ridge Feature Ablation (Mean $R^2$ Delta Across 6 Thresholds)}
\label{tab:featureablation}
\centering
\begin{tabular}{lcl}
\toprule
Dropped Feature & Mean $R^2$ $\Delta$ & Most-Affected Threshold \\
\midrule
meanNameLength & $+0.053$ & prefixSimilarityMinShared \\
prefixSimilarityCoeff & $+0.020$ & compressMinLen \\
repeatRate & $+0.015$ & prefixSimilarityMinShared \\
scopeDepthVariance & $-0.005$ & inlineMaxLen \\
nameEntropy & $+0.008$ & prefixSimilarityMinShared \\
\bottomrule
\end{tabular}
\end{table}

\textbf{Honest reading.} A positive $\Delta$ means dropping that feature
\emph{improved} test-set $R^2$ on average -- i.e., the feature was net
noise for that threshold on this 300-workload training set, not a useful
signal. Every feature except \texttt{scopeDepthVariance} shows a positive
mean delta, meaning the full 5-feature model is, on average, very slightly
outperformed by simpler 4-feature ablations. \texttt{scopeDepthVariance} is
the only feature with a negative mean delta ($-0.005$, i.e. removing it
hurts, if only marginally), concentrated in \texttt{inlineMaxLen}
($-0.027$), making it the one feature with a defensible claim to adding
signal beyond noise. This is a modest result, not a strong one: at 300
training workloads and 5 already-correlated features (e.g.
\texttt{prefixSimilarityCoeff} and \texttt{meanNameLength} both partly
encode "how compressible is this workload"), the ablation deltas are small
relative to the R^2 values themselves (Table~\ref{tab:r2}) and should be
read as a first-pass sensitivity check, not a definitive feature-importance
ranking -- more training workloads or a regularization-path analysis would
be needed to say more confidently which features are load-bearing.

\subsection{Review-3 Extension: Correctness Fix --- Lookup Cache and Scope Shadowing}
\textbf{Motivation.} Reviewing Review-2's \texttt{LRULookupCache} before
extending it (Section VIII-I below) surfaced a correctness question that
Review-2's own tests did not exercise: the cache is keyed by
\texttt{(hash, name)} with no representation of scope, so it is not obvious
from inspection alone that a symbol shadowing an outer-scope declaration of
the same name is resolved correctly.

\textbf{Bug found.} It is not resolved correctly. \texttt{insert()} never
touched \texttt{lookupCache\_} prior to this revision; only
\texttt{exitScope()} invalidated cache entries, and only for symbols in the
scope actually being popped. A new declaration that shadows an
already-cached outer symbol of the same name therefore leaves the outer
symbol's stale \texttt{(hash, name) -> id} mapping in place, and
\texttt{lookup()}'s cache fast path returns it instead of resolving through
the scope stack. Confirmed empirically rather than by inspection alone: a
targeted reproduction forces an outer declaration of \texttt{"x"} to a
promotable (COMPRESSED) representation and a shadowing inner \texttt{"x"}
to a non-promotable one (INTERNED, via the exact-repeat rule, which takes
priority over every other branch of \texttt{decide()}), then issues three
\texttt{lookup("x")} calls from inside the inner scope. Before the fix,
\texttt{promotions()} reports 1 -- the three accesses were credited to the
shadowed \emph{outer} entry, which is only reachable at all if the cache
served its stale mapping instead of the live inner symbol. After the fix,
\texttt{promotions()} reports 0, matching correct scope-respecting
resolution (the inner symbol is already INTERNED and has nothing to
promote to).

\textbf{Fix.} \texttt{insert()} now invalidates \texttt{lookupCache\_}'s
entry for a hash immediately after computing it, before the new entry is
added to the scope's index -- so any declaration that reuses a hash,
shadowing or otherwise, forces the next lookup for that hash to re-resolve
through the real scope stack rather than serve a cached answer that predates
the new declaration. A regression test
(\texttt{test\_lookup\_cache\_respects\_scope\_shadowing}) encodes the
reproduction above and is now part of the permanent suite. This is reported
as a correctness fix, not a performance change: it does not alter any
memory or compression figure in this paper, and the cache's throughput
characteristics (Section VIII-E) are unaffected, since invalidating one
hash on every insert is $O(1)$.

\textbf{Why Review-2's own tests missed it.} The existing cache-invalidation
test uses a symbol name that exists only inside the scope under test; it
never declares the same name in two enclosing scopes. This is stated
plainly because it is a useful, general lesson independent of this specific
bug: a passing test suite supports the claim ``the mechanism is tested,''
not the stronger claim ``the mechanism is tested for every interaction that
matters,'' and the gap between those two claims is exactly where this bug
lived.

\subsection{Review-3 Extension: Hashing and Storage Algorithm Comparison}
\textbf{Motivation.} Review-2's Future Work (Section XV) proposed
``reducing cold-lookup cost directly (e.g., a secondary index keyed by a
cheap partial reconstruction) rather than only amortizing it via caching''
as an alternative to the lookup cache. This revision implements and
measures two concrete structural alternatives, plus isolates the index hash
function itself as a variable BUDGET-SYM had not previously varied.

\textbf{Hash function.} \texttt{BudgetSym} is templated on its index hash
function (\texttt{BudgetSymT<HashFn>}, defaulting to the original FNV-1a via
a type alias, so no existing caller's behavior changes). Two additional
backends are compared: MurmurHash3 \cite{appleby-murmur} (the public-domain
32-bit x86 finalizer, applied twice with different seeds and packed into a
64-bit key, since implementing the full MurmurHash3\_x64\_128 variant was
judged out of proportion to what this lookup path needs) and DJB2, a
single-multiply-and-add hash with no mixing step, included as a
cheap-end comparison point.

\textbf{Robin Hood open addressing.} \texttt{RobinHoodSymbolTable}
\cite{celis-robinhood} is a new, direct string-keyed baseline (no adaptive
representation selection, no interning) applying Robin Hood linear-probing
open addressing the same way \texttt{ConventionalSymbolTable} applies
\texttt{std::unordered\_map} -- isolating table \emph{structure} as the only
variable against the conventional baseline. Entries live inline in one
contiguous array rather than one heap node per chain-bucket entry, which
is the actual mechanism behind its lookup speed, not a faster hash. This
comes at a real, separately-reported memory cost: every slot, occupied or
not, is a live \texttt{std::string} plus metadata in a preallocated array,
so the structure carries genuine ``slack'' memory for its empty slots that
a chaining table's null-pointer buckets do not pay. Deletion is not
implemented (no tombstones, no backward-shift delete), since this project's
only removal path is whole-scope reclamation, identical to
\texttt{ConventionalSymbolTable}'s own approach -- there is no exercised
code path that would call a single-key delete from a still-live table.

\textbf{Shared trie.} \texttt{TrieSymbolTable} \cite{fredkin-trie} shares
identifier prefixes across every string ever inserted, project-wide,
independent of insertion order -- strictly more powerful sharing than
BUDGET-SYM's COMPRESSED front-coding, which only shares a prefix with the
one identifier inserted immediately before it in the chain. The cost is
real: a dedicated node with its own child map is charged per distinct
branching character, which is more expensive than front-coding's raw
suffix bytes when there is little global prefix overlap to exploit. Lookup
needs no reconstruct-then-compare step at all -- walking the query
string's characters down existing children \emph{is} the equality check --
which is a structurally different answer to the lookup-cost question than
either the LRU cache (Section VIII-E) or a faster hash function: it avoids
reconstruction entirely, for every lookup, hit or miss, rather than
amortizing it across repeats of the same symbol.

\begin{table}[htbp]
\caption{Best Algorithm per Dataset, by Metric}
\label{tab:algocompare}
\centering
\small
\begin{tabular}{lcc}
\toprule
Dataset & Best Compression & Best Lookup (hit) \\
\midrule
small & BudgetSym, 1.65$\times$ & RobinHood, 0.054 $\mu$s \\
medium & BudgetSym, 1.63$\times$ & RobinHood, 0.055 $\mu$s \\
large & BudgetSym, 1.64$\times$ & RobinHood, 0.054 $\mu$s \\
high-prefix-sim. & BudgetSym, 2.37$\times$ & RobinHood, 0.082 $\mu$s \\
random-identifiers & BudgetSym, 1.34$\times$ & RobinHood, 0.061 $\mu$s \\
nested-scopes & BudgetSym, 1.84$\times$ & RobinHood, 0.054 $\mu$s \\
hot-cold-access & BudgetSym, 1.57$\times$ & RobinHood, 0.046 $\mu$s \\
memory-stress & BudgetSym, 1.38$\times$ & RobinHood, 0.105 $\mu$s \\
\bottomrule
\end{tabular}
\end{table}

\textbf{Results.} \texttt{src/algorithm\_benchmark\_main.cpp} compares
BudgetSym-FNV1a, BudgetSym-Murmur3, BudgetSym-DJB2, RobinHood, and Trie on
the same eight datasets and cost-accounting model as
Table~\ref{tab:cache}. Table~\ref{tab:algocompare} reports the best
algorithm per dataset on each of the two headline metrics. RobinHood wins
lookup latency on every dataset, typically by roughly an order of magnitude
or more over any BudgetSym variant, and BudgetSym wins compression on every
dataset; no algorithm in the comparison wins both on any dataset. This is
the direct, structural consequence of what each design optimizes for, not
a coincidence of these particular measurements: BudgetSym pays a
reconstruction cost specifically because it stores compressed or shared
representations a direct string key cannot; RobinHood pays no
reconstruction cost specifically because it stores full, uncompressed
strings with no sharing, and pays real slack-memory overhead for that
instead. Trie wins neither metric outright on any dataset but rejects
absent identifiers fastest of all five algorithms on most datasets, since a
missing child early in the walk ends the search immediately.

\textbf{On hash function choice.} Memory is bit-for-bit identical across
all three BudgetSym hash variants on every dataset and every run, as
expected (the hash affects only bucket placement, never what is stored).
Which hash function is fastest on lookup changes on every rerun of this
benchmark on the shared, non-isolated development machine used for this
work; the ranking itself is not reproducible, which is reported as the
finding rather than papered over by presenting one run's ranking as
definitive. The interpretation offered is that identifier strings in this
project's workloads are short enough, and the dominant lookup cost
(reconstruction) large enough, that hash-computation cost is not a
distinguishing factor at this data scale.

\textbf{An overclaim caught before publication.} An early run of this
comparison showed Trie's lookup latency on the \texttt{random-identifiers}
dataset higher than all three BudgetSym variants, consistent with a
plausible mechanistic story: BudgetSym's adaptive policy can fall back to a
cheap representation for identifiers that do not compress well, while a
trie has no equivalent escape hatch. Repeating the identical run three more
times showed Trie faster than BudgetSym twice and slower once, on identical
code and data -- the effect was single-run timing noise on a shared
machine, not a real, reproducible finding, and is reported here as such
rather than as the clean story the first run suggested. This is offered as
a concrete instance of the general risk Threats to Validity item 1 already
names (single-seed measurement): a plausible-sounding, mechanism-backed
result is not automatically a true one.

\section{Experimental Methodology}
Two real baselines are used: a conventional hash table
(\texttt{unordered\_map<string, Symbol>}, one full string copy per symbol,
used as the 1.00$\times$ reference point) and a string-interning table.
Eight synthetic benchmark datasets are used: small, medium, large,
high-prefix-similarity, random-identifiers, nested-scopes, hot-cold-access,
and memory-stress, each designed to isolate a different identifier
characteristic. Review-2 added three real-world codebase datasets and
extended from single-seed to 30-seed evaluation with confidence intervals.
This revision adds a ninth dataset, ml-predicted, and a fourth
implementation variant, BudgetSym-WithCache, both described in Section VIII.

For every experiment, the independent variable is the symbol-table
implementation; the dependent variables are tracked memory and lookup
latency; and the identifier workload, operation sequence, and
memory-accounting model are held constant across implementations.

\section{Evaluation Metrics}
Memory metrics: total tracked memory, tracked memory per symbol, and
compression ratio relative to the conventional baseline. Performance
metrics: lookup latency per dataset, and lookup-cache hit rate
[\textbf{new}]. Statistical metrics: mean compression ratio across 30 seeds,
95\% confidence intervals, and $p$-values. Adaptive metrics: number of
representation promotions, reclaimed memory on scope exit, memoized-
reconstruction hit/cold-lookup counts [\textbf{new}], and predicted-versus-
hand-picked/grid-search compression ratio [\textbf{new}].

\section{Results}
\subsection{Memory: BudgetSym vs. Two Baselines}
BudgetSym reduces tracked memory relative to both baselines on every
dataset tested, with the largest gains on prefix-similar and nested-scope
workloads, consistent across Review-1, Review-2, and this revision (memory
accounting is unaffected by the cache/profiler/predictor additions --- see
Section VII).

\subsection{Lookup Latency}
See Table~\ref{tab:cache} (Section VIII-E). BudgetSym's lookup latency
remains higher than the conventional baseline's on every dataset measured,
but the 64-entry lookup cache measurably reduces it relative to BudgetSym
without the cache, on a hot-path access pattern.

\section{Ablation Study}
Table~\ref{tab:ablation} reports a four-variant ablation on a fixed
workload. Disabling scope-aware reclamation increases tracked memory
substantially --- the largest single-mechanism effect, confirming
scope-based reclamation as the dominant contributor. Disabling
access-frequency-driven promotion changes memory only marginally, since
promotion's direct memory cost is small; its value is primarily in lookup
latency for hot symbols. Removing adaptive selection entirely is the
worst-performing configuration. This revision's \texttt{disableMLThreshold-
Prediction} flag (added to \texttt{PolicyConfig}) keeps the
NoAccessFrequency variant's fixed threshold from being silently overwritten
by the automatic post-100-symbol predictor described in Section VIII-F ---
without it, this ablation variant's isolation of the access-frequency
mechanism would be broken by exactly the new mechanism this revision adds.

\begin{table}[htbp]
\caption{Ablation Results (Fixed Workload, R1 Thresholds)}
\label{tab:ablation}
\centering
\begin{tabular}{lcc}
\toprule
Variant & Tracked Mem. & Promotions \\
\midrule
BudgetSym-Full & 31{,}894 B & 4 \\
BudgetSym-NoScope & 46{,}496 B & 4 \\
BudgetSym-NoAccessFrequency & 34{,}143 B & 0 \\
BudgetSym-NoAdaptiveSelection & 57{,}042 B & 0 \\
\bottomrule
\end{tabular}
\end{table}

\section{Discussion}
The results indicate a genuine memory--latency trade-off rather than an
unconditional win. BUDGET-SYM reduces tracked memory relative to both
baselines on every dataset tested. Review-2's threshold optimization
demonstrated that the hand-picked constants in Review-1 left measurable
performance on the table. Real-world evaluation on FreeRTOS, Arduino, and
Zephyr confirmed that the gains observed on synthetic data transfer to
actual embedded codebases. Multi-seed validation established statistical
significance across all eight synthetic datasets.

This revision's two additions attack the lookup-latency side of the
trade-off directly rather than accepting it as fixed. The lookup cache is a
genuine, measured latency reduction, though it is a mitigation, not a
resolution: BudgetSym's lookup latency still exceeds the conventional
baseline's on every dataset tested. The threshold predictor demonstrates
that a workload's own profile, sampled online in the first 100 symbols, can
recover most of a dataset-specific grid search's compression benefit without
ever running that grid search against the real workload. Its main
limitation is symmetrical to the grid search's: the grid search overfits to
eight fixed datasets, and the predictor's training data is itself 300
synthetic workloads, not real compiler traffic --- see Threats to Validity,
item 8.

\section{Threats to Validity / Limitations}
\begin{enumerate}
\item Real-world workloads evaluated with a single seed in Review-2; not
  re-run for this revision. Multi-seed validation for real corpora remains
  future work.
\item Tracked memory vs. OS RSS. The memory model is a deterministic
  accounting construct, not OS resident memory, and this revision's cache
  and memoization additions are not charged against it (see Section VII).
\item Lookup overhead is reduced, not resolved. The 64-entry LRU cache
  measurably lowers BudgetSym's own lookup latency on a hot-path access
  pattern (Section VIII-E), but BudgetSym's lookup latency still exceeds the
  conventional baseline's on every dataset measured; the cache mitigates the
  structural cost, it does not eliminate it.
\item COMPRESSED chain-interior reclamation is incomplete. Interior nodes
  accumulate until the next re-anchor point (unchanged from Review-1/2).
\item Prototype, not production integration. BUDGET-SYM has not been
  integrated into a real compiler front end.
\item Prior-art search incomplete. A full IEEE Xplore/ACM DL/arXiv search
  has not yet been completed.
\item Grid search (Section VIII-A) is tuned on eight fixed synthetic
  datasets; generalization to out-of-distribution workloads beyond what the
  threshold predictor addresses is untested.
\item ML threshold predictor trained on synthetic workloads only. The 300
  training workloads (Section VIII-F) are themselves synthetically generated
  from randomized feature targets, using a reduced 144-combination grid
  search as the ground-truth labeler rather than the full 15,000-combination
  sweep. Generalization to real codebases is reported on one dataset
  (Table~\ref{tab:mlcompare}) but is not exhaustively validated against the
  FreeRTOS/Arduino/Zephyr corpora or against multiple random seeds.
\item A latent benchmark-harness measurement bug (Section VIII-E) means
  lookup-latency numbers from Review-1 and Review-2 were not reliable; only
  this revision's own lookup-latency figures should be trusted.
\item A correctness bug in the lookup cache (Section VIII-G) meant a symbol
  shadowing an outer-scope declaration of the same name could resolve to
  the wrong, stale entry; fixed and regression-tested in this revision, but
  it shipped in Review-2 undetected, which is itself evidence that this
  project's test coverage claims should be read as ``tested for the cases
  written,'' not ``tested exhaustively.''
\item The grid-search-based ablation study (Section XI) and the ML
  threshold predictor both interact with \texttt{PolicyConfig} at runtime,
  but \texttt{ablation\_main.cpp}'s \texttt{BudgetSym-Full} and
  \texttt{BudgetSym-NoScope} variants do not set
  \texttt{disableMLThresholdPrediction}, while \texttt{BudgetSym-NoAccessFrequency}
  explicitly does (with a comment explaining why). The asymmetry means two
  of the four ablation variants can have their thresholds silently replaced
  by the ML predictor mid-run while the other two cannot, which is a
  potential confound in that comparison not yet resolved.
\item Robin Hood's ``slack'' memory (Section VIII-H) is reported as a
  separate diagnostic, not folded into the shared tracked-memory
  accounting model every other implementation in this paper uses, so it is
  not directly comparable to the memory-per-symbol figures in
  Table~\ref{tab:algocompare}'s companion CSV without consulting that
  separate figure.
\end{enumerate}

\section{Future Work}
Planned next steps include:
\begin{itemize}
\item Multi-seed statistical validation of real-codebase results, and of
  this revision's cache-hit-rate and ML-predicted-compression-ratio claims.
\item Charging the lookup cache's and memoization cache's own memory cost
  against the tracked-memory model, so the latency/memory trade-off this
  revision introduces is fully accounted for rather than latency-only.
\item Training the ThresholdPredictor directly on real embedded-corpus
  identifier traffic (FreeRTOS/Arduino/Zephyr) instead of purely synthetic
  workloads, and re-running the full 15,000-combination grid search per
  training workload if compute allows.
\item Fix COMPRESSED chain-interior reclamation via a compacting GC pass
  over \texttt{compPool\_}.
\item \textbf{Addressed this revision:} reducing cold-lookup cost directly
  rather than only amortizing it via caching. RobinHood (Section VIII-H)
  eliminates the reconstruct-and-compare step entirely by giving up
  compression, quantifying exactly what that trade costs and buys; the
  originally-proposed narrower idea (a secondary index keyed by a cheap
  partial reconstruction, kept compressed) remains open and would need to
  close the gap without RobinHood's full slack-memory cost.
\item Resolving the ablation/ML-threshold-predictor interaction flagged in
  Threats to Validity so all four ablation variants are compared under a
  consistently fixed or consistently adaptive policy.
\item Reducing the corpus fixtures' repository footprint (the extracted
  identifier lists are currently committed directly rather than
  regenerated from the already-gitignored source corpora on demand).
\item Completed IEEE Xplore/ACM DL/arXiv literature and prior-art search.
\item Integration with a real compiler front end (e.g., as an LLVM pass
  component).
\end{itemize}

\section{Conclusion}
This paper reports BUDGET-SYM, a prototype adaptive, scope-aware
symbol-table architecture for memory-constrained embedded compilers.
Measured against two real baselines on eight synthetic benchmark datasets,
the prototype reduces tracked memory by 1.24$\times$--2.59$\times$. Review-2
contributed a systematic grid-search threshold optimization, evaluation
against real embedded C corpora, and multi-seed statistical validation. This
revision contributes two further advances: a 64-entry LRU lookup cache
(with per-entry memoized COMPRESSED reconstruction) that measurably reduces
BudgetSym's own lookup latency on hot-path access patterns without changing
its memory-accounting model, and an online WorkloadProfiler feeding an
offline-trained Ridge-regression ThresholdPredictor that automatically
reconfigures policy thresholds after a 100-symbol sampling window, recovering
most of a dataset-specific grid search's compression benefit on an unseen
workload with no runtime ML dependency. This revision (Review-3) contributes
a correctness fix -- a scope-shadowing bug in the Review-2 lookup cache,
found by review and confirmed empirically before being fixed and
regression-tested -- and a five-way comparison of alternative hashing and
storage backends (two additional index hash functions, a Robin Hood
open-addressing baseline, and a shared-trie baseline) directly answering
Review-2's own proposed future work on reducing cold-lookup cost
structurally rather than only amortizing it. The comparison's central
finding is that no algorithm wins both memory and lookup latency on any of
the eight datasets: Robin Hood wins lookup latency universally by forgoing
compression, and BudgetSym wins compression universally by paying the
reconstruction cost that Robin Hood avoids. Lookup latency overhead is
reduced along two independent axes (caching, and now this structural
comparison) but remains an honestly reported, unresolved cost relative to
the conventional baseline for BudgetSym itself. Further work toward LLVM
integration, real-corpus predictor training, resolving the
ablation/ML-predictor interaction noted in Threats to Validity, and a
complete prior-art search are the required next steps.

\begin{thebibliography}{00}
\bibitem{liu-lab2} Link\"oping University, ``Lab 2: Symbol Table --- TDDE66
  Compiler Construction,'' compiler-construction documentation. [Online].
  Available: https://www.ida.liu.se/\textasciitilde{}TDDE66/laboratories/instructions/lab2.html
\bibitem{jones-notes} D. W. Jones, ``Compiler Construction --- The Symbol
  Table,'' University of Iowa, compiler notes. [Online]. Available:
  https://homepage.cs.uiowa.edu/\textasciitilde{}dwjones/compiler/notes/09.shtml
\bibitem{rust-interning} Rust compiler ecosystem, ``Interning / Symbol
  Table,'' documentation and implementation discussion. [Online]. Available:
  https://patterns.totorojam.com/patterns/interning/
\bibitem{gcc-msp430} GNU Project, ``MSP430 Options,'' GCC documentation,
  memory-constrained embedded-target options. [Online]. Available:
  https://gcc.gnu.org/onlinedocs/gcc-10.4.0/gcc/MSP430-Options.html
\bibitem{aho} A. V. Aho, M. S. Lam, R. Sethi, and J. D. Ullman,
  \emph{Compilers: Principles, Techniques, and Tools}, 2nd ed. Pearson, 2006.
\bibitem{ieee-structure} IEEE Author Center Conferences, ``Structure Your
  Paper.'' [Online]. Available:
  https://conferences.ieeeauthorcenter.ieee.org/write-your-paper/structure-your-paper/
\bibitem{ieee-tools} IEEE Author Center Conferences, ``Authoring Tools and
  Templates.'' [Online]. Available:
  https://conferences.ieeeauthorcenter.ieee.org/write-your-paper/authoring-tools-and-templates/
\bibitem{celis-robinhood} P. Celis, P.-A. Larson, and J. I. Munro, ``Robin
  Hood Hashing,'' in \emph{Proc. 26th Annual Symposium on Foundations of
  Computer Science (FOCS)}, 1985, pp. 281--288.
\bibitem{fredkin-trie} E. Fredkin, ``Trie Memory,'' \emph{Communications of
  the ACM}, vol. 3, no. 9, pp. 490--499, 1960.
\bibitem{appleby-murmur} A. Appleby, ``MurmurHash3 and the SMHasher hash
  function test suite.'' [Online]. Available:
  https://github.com/aappleby/smhasher
\end{thebibliography}

\vspace{1em}
\noindent\textit{Formatting note: this document is an IEEE-inspired academic
draft prepared for internal faculty review. It is not an official IEEE
publication, has not been submitted to or accepted by any IEEE venue, and
should not be represented as such.}

\end{document}
